\chapter{A Consensus Protocol Skeleton}\label{app:protocol}

This appendix sketches the mechanism that \S\ref{sec:mechanism} argues is required and
\S\ref{sec:consensus} models: the procedure by which a group of independent parties converts an
intent into a specification that others may then execute against. It is written because the
chapter's central claim---that the intent layer's guarantee comes from adversarial consensus rather
than from a verifier---is hollow without at least a sketch of what such a procedure looks like, and
because the claim that the intersection is \emph{unoccupied} (Chapter \ref{ch:literature}) is
easier to assess against a concrete proposal than against a description.

Three things this appendix is not. It is not an implementation; no part of it has been built or
run. It is not validated; no experiment is reported, and the parameters are stated as quantities to
be fixed rather than as values that were measured. And it is not a claim about correctness: the
protocol below can be followed exactly and still produce a specification that misrepresents the
intent, for the reasons Proposition \ref{prop:three-grounds} sets out. What it does claim is a
shape: that the obligations the rest of the paper places on the intent layer can be discharged by a
procedure with these steps, and that each step has an analogue in deployed practice.

\section{The Object Being Attested}

The protocol produces an \term{attestation record}, which is the intent-layer artefact the rest of
the chain consumes. Its fields are the intent-layer counterpart of the artefact evidence list
(Table \ref{tab:evidence-package}), and the correspondence is deliberate: the two lists are read by
the same downstream machinery.

\begin{center}
\small
\begin{longtable}{L{0.9cm}L{5.2cm}L{8.5cm}}
\caption{Attestation record}\label{tab:attestation}\\
\toprule
\textbf{\#} & \textbf{Field} & \textbf{Content and why it is required} \\
\midrule
\endfirsthead
\toprule
\textbf{\#} & \textbf{Field} & \textbf{Content and why it is required} \\
\midrule
\endhead
\bottomrule
\endlastfoot
1 & Specification text & The formal object $\Phi$ itself, with its scope $\Sigma$ and its declared
  environmental assumptions $A$. \\
2 & Intent statement & The requirement as stated by the party with standing, in whatever notation that
  party used. Kept verbatim: paraphrase is where intent is lost, and the gap between fields 1 and 2
  is precisely the object Proposition \ref{prop:three-grounds}(i) says cannot be closed
  mechanically. \\
3 & Attestor identities and stakeholder classes & Who attested, and which interest each
  represents. Required because the reliability of the aggregate depends on the \emph{independence}
  of the error sources, and stakeholder class is the observable proxy for that independence
  (\S\ref{sec:mechanism}). \\
4 & Evidence provenance & For each attestor, the sources on which the judgement rested. This is the
  field that makes $q$ measurable in principle: two attestors relying on the same source are
  correlated whether or not they know it, which is the common-cause mechanism of
  \cite{dietrich2004}. \\
5 & Dissent record & Objections raised and how they were resolved. A record with no dissent is
  evidence about the procedure, not about the specification. \\
6 & Threshold and rule & The decision rule used, the threshold, and its rationale. \\
7 & Authorising body & Who adopted the specification and with what standing
  (\S\ref{sec:authority}). Distinct from the attestors: attesting is judging, authorising is
  deciding. \\
8 & Validity horizon and re-confirmation triggers & When the attestation lapses, and which events
  force re-attestation rather than mere re-verification (\eqref{eq:spec-decay}). \\
9 & \textbf{Known unresolved intent} & Requirements the procedure could not settle. The analogue of
  uncovered behaviour in Table \ref{tab:evidence-package}, item 9, and the field most likely to be
  omitted because recording it weakens the appearance of the result. \\
10 & \textbf{Re-attestation rules} & Which changes to the specification, the intent, the
  environment or the authorising body invalidate the record. \\
\end{longtable}
\end{center}

Fields 9 and 10 carry the same weight here as items 9 and 10 carry there. A record that omits field
9 is not a stronger record; it is a record whose boundary is unknown, and the downstream evidence
chain will treat the specification as complete when it is not.

\section{The Procedure}

The steps below are stated in the order in which they are performed. Steps 1--3 are constructive,
step 4 is adversarial, step 5 is the decision, and steps 6--7 are maintenance.

\begin{enumerate}[label=\textbf{P\arabic*.},leftmargin=*]
\item \term{Frame the intent.} The party with standing states the goal, the constraints and the
      acceptance criteria, and names the scope $\Sigma$. Output: field 2.
\item \term{Derive the specification.} One or more proposers produce candidate specifications from
      the intent. This step may be automated, and in the framework of Chapter \ref{ch:program} it
      is where generation is expected to be cheapest. Output: candidate values of field 1.
\item \term{Make the derivation traceable.} Each clause of the candidate specification is linked to
      the part of the intent statement it discharges. Clauses with no link are unrequested; parts
      of the intent with no clause are unrepresented. This step is the mechanism by which field 9
      becomes writable rather than aspirational, and it is the discipline that recent
      traceable-refinement work automates (\S\ref{sec:gaps}).
\item \term{Attest adversarially.} Attestors drawn from deliberately different stakeholder classes
      judge the candidate against the intent, with an explicit instruction to find failures rather
      than to confirm. Output: fields 3, 4 and 5. Two design choices matter more than the rest.
      Attestors must be selected for \emph{diversity of interest and of evidence}, because that is
      what the reliability argument requires, not merely for competence; and dissent must be
      recorded rather than resolved away, because the record of an unresolved disagreement is the
      most informative output of the step.
\item \term{Decide and authorise.} A threshold rule is applied to the attestations, and the
      authorising body adopts the specification. The two acts are distinct and must be separately
      recorded (fields 6 and 7), for the reason \S\ref{sec:authority} gives: a specification can be
      correct and unauthorised, or authorised and wrong.
\item \term{Set the horizon and the triggers.} The record is given a validity horizon and a list of
      events that force re-attestation. The horizon is not cosmetic; eq.~\eqref{eq:spec-decay}
      makes the assurance condition time-indexed precisely because an attested specification can
      drift from a moving intent while every mechanical check continues to pass.
\item \term{Re-attest on trigger.} On any trigger from field 10, the procedure restarts at step 1
      for the affected clauses. Re-verification of the implementation is not a substitute: if the
      intent moved, the artefact that was verified is no longer the right artefact to verify.
\end{enumerate}

\section{What Each Step Buys, and What It Does Not}

The value of writing the protocol out is that it separates the parts of the problem that different
instruments actually address, which the chapter's argument needs and a prose description obscures.

\emph{Steps 2 and 3 are largely mechanisable}, and are where automation has the most to offer. They
address the expressibility question. Nothing in this paper denies that they can be automated well;
the framework expects it.

\emph{Step 4 is not mechanisable and cannot be made so}, because its input includes judgements about
what the intent should have been. This is the discovery question, and the protocol's contribution to
it is procedural: it forces the disagreement to be recorded and the evidence provenance to be
declared, which converts an unmeasurable quantity into an estimable one. That is the most that can
be claimed, and it is less than the reliability argument of \S\ref{sec:consensus} would need to be
a prediction.

\emph{Step 5 is a different kind of act from steps 1--4}, and this is the point \S\ref{sec:authority}
adds to the framework. No amount of correct judgement by appropriately diverse attestors makes the
result binding; standing is not earned by accuracy. A protocol that recorded attestations but not
authorisation would have no answer to the objection that a specification was never anyone's to
adopt.

\emph{Steps 6 and 7 are the parts most often absent from deployed practice}, and they are the parts
that a lifecycle view makes unavoidable. Their absence is invisible at the moment of attestation and
apparent only afterwards, when the intent has moved and the record still says what it said.

\section{Failure Modes of the Protocol Itself}

A protocol sketch that listed only its steps would be a sales document. The modes below are the ones
this design is most exposed to, stated as properties of the procedure rather than of its
participants.

\begin{enumerate}[label=(\alph*)]
\item \term{Selection capture.} If the set of attestors is chosen by the party whose specification
      is under judgement, diversity can be satisfied nominally while the errors remain correlated.
      Field 3 does not detect this, because a captured roster is still a roster. The mitigation is
      procedural and outside the protocol: the selection rule must be fixed before the candidate is
      known.
\item \term{Provenance collapse.} Attestors from different stakeholder classes become dependent if
      they draw on the same source, and field 4 is what makes this visible. Where provenance is
      undeclared the protocol degenerates to the Condorcet case, whose premises are not jointly
      satisfiable in general \cite{dietrich2008}.
\item \term{Threshold theatre.} A high threshold over a small, correlated roster yields a
      confidence that the bound of \eqref{eq:chernoff} does not support, because the bound's
      effective sample size is set by $\bar\rho$ and not by the count. This is the design-effect
      point of Remark \ref{rmk:design-effect} appearing as a governance failure.
\item \term{Unfalsifiable acceptance.} If no observable consequence is specified at step 1, the
      attestation cannot later be scored, and the protocol produces a record that is
      indistinguishable from a correct one. This is the limitation Remark
      \ref{rmk:consensus-status} records for the whole consensus layer, and it is the strongest
      argument against relying on the protocol in domains that lack scoreable outcomes.
\item \term{Silent clause loss.} Step 3 makes unrepresented intent visible only if the intent
      statement is itself complete. Where the party with standing cannot say what it wants until it
      sees what it does not want, field 9 is written from a moving target, and the record's
      completeness claim is weaker than its format suggests.
\end{enumerate}

None of these is repaired by a better threshold. Counting them: two are properties of the roster and
its provenance, one is a property of the decision rule, and two are properties of the domain. That
distribution is the honest summary of what a consensus protocol can and cannot contribute: it can
make the relevant properties \emph{recorded and therefore auditable}, and it cannot make them
\emph{true}.
